\documentclass[11pt]{article}
\usepackage[paperwidth=220mm,paperheight=320mm,left=16mm,right=16mm,top=18mm,bottom=18mm,headheight=16pt,headsep=10pt,footskip=14pt]{geometry}
\usepackage{fontspec,graphicx,tikz}
\usepackage[unicode,hidelinks]{hyperref}
\usepackage{polyglossia}
\setmainlanguage{latin}
\setotherlanguages{arabic}
\setmainfont{Linux Libertine O}[Ligatures=TeX]
\newfontfamily\arabicfont{Amiri}[Script=Arabic]
\newfontfamily\ArBody{Amiri}[Script=Arabic]
\newfontfamily\OrnFont{FreeSerif}
\newfontfamily\NallinoSigns{NallinoSigns.otf}[Path=./fonts/]
\newcommand{\AbjadZero}{{\NallinoSigns\symbol{"E000}}}
\hypersetup{pdftitle={Al-Battani, Kitab az-Zij as-Sabi'. Textus Arabicus (Nallino, Pars III). S11},pdfauthor={Al-Battani; Carlo Alfonso Nallino (ed.)}}
\makeatletter
\def\ps@sourceedition{%
\def\@oddhead{\vbox{\hbox to\textwidth{\fontsize{8}{10}\selectfont AL-BATTĀNĪ — TEXTUS ARABICUS (NALLINO, PARS III)\hfil\leftmark}\vskip3pt\hrule height.25pt}}%
\let\@evenhead\@oddhead
\def\@oddfoot{\hbox to\textwidth{\fontsize{7}{9}\selectfont\rightmark\hfil\thepage}}%
\let\@evenfoot\@oddfoot}
\makeatother
\pagestyle{sourceedition}
\setlength{\parindent}{0pt}\setlength{\parskip}{0pt}
\tracinglostchars=2
\newcommand{\DSource}[2]{\clearpage\markboth{#2}{AB01-PDF#1}\pdfbookmark[0]{AB01-PDF#1 — #2}{bk-#1}\null}
% place #3 with its anchor #1 at (#2) from the top left corner of the page
\newcommand{\At}[3]{\begin{tikzpicture}[remember picture,overlay]\node[anchor=#1,inner sep=0pt,outer sep=0pt] at ([shift={(#2)}]current page.north west) {#3};\end{tikzpicture}}
\newsavebox{\LBox}
% a right-to-left line #2 justified to the width #1 (\RLJ) or at its natural width (\RLN); narrowed to #1 if wider
\newcommand{\RLJ}[2]{\sbox{\LBox}{\RL{#2}}%
 \ifdim\wd\LBox>#1\relax\typeout{NARROWED|\the\wd\LBox|#1}\resizebox{#1}{\ht\LBox}{\usebox{\LBox}}%
 \else\makebox[#1][s]{\RL{#2}}\fi}
\newcommand{\RLN}[2]{\sbox{\LBox}{\RL{#2}}%
 \ifdim\wd\LBox>#1\relax\typeout{NARROWED|\the\wd\LBox|#1}\resizebox{#1}{\ht\LBox}{\usebox{\LBox}}%
 \else\usebox{\LBox}\fi}
% a letter-spaced name in the notes (R e i n a u d on p. 25: 2.55 pt between the letters at 9.2 pt)
\newcommand{\Name}[1]{{\addfontfeatures{LetterSpace=27}#1}}
\newcommand{\Note}[1]{\textsuperscript{\fontsize{8}{8}\selectfont\addfontfeatures{Numbers=Lining}#1}}
\newcommand{\Ov}[1]{\vbox{\hrule height .4pt\kern 1.0pt\hbox{#1}}}
\newcommand{\ArRun}[1]{\textarabic{#1}}
% a display line in Latin type #3 scaled to the width #1 and the height #2 of its printed line
\newcommand{\LFit}[3]{\sbox{\LBox}{#3}\resizebox{#1}{#2}{\usebox{\LBox}}}
\begin{document}
\thispagestyle{empty}
\begin{center}\Large AL-BATTĀNĪ\par\vspace{4mm}\LARGE KITĀB AZ-ZĪJ AṢ-ṢĀBIʾ\par
\vspace{5mm}\large textus Arabicus a Carolo Alphonso Nallino editus (pars III)\par\vspace{8mm}
\Large S11 --- editio diplomatica\par\vspace{4mm}\large Paginae impressae 80--81\end{center}
\vspace{12mm}\noindent Singulae lineae paginarum impressarum singulis lineis huius editionis respondent et ibi
collocantur ubi in pagina impressa stant (lineae quae totam latitudinem implent ad eandem latitudinem
extenduntur). Textus Arabicus cum signis vocalium, ut impressa sunt, ex imaginibus paginarum transcriptus est;
productiones litterarum (kashida) non transcribuntur. Numeri foliorum codicis et numeri linearum Nallini in
marginibus stant ubi impressi sunt. Haec transcriptio a nullo homine recognita est.\par
\DSource{1072}{PAG. 80}
\At{base east}{499.33pt,-141.09pt}{\hypertarget{AB01-PDF1072-L01}{}{\fontsize{13.8}{27}\ArBody\RLJ{420.10pt}{جزءًا فبه يكون الخمسة الاجزاء والرُّبْع\Note{1} تسعًا وثلثين درجة واثنتين وعشرين دقيقة وذلك هو مقدار}}}
\At{base east}{499.33pt,-168.09pt}{\hypertarget{AB01-PDF1072-L02}{}{\fontsize{13.8}{27}\ArBody\RLJ{420.10pt}{خطّ \Ov{ه ف} الذي من مركز الارض الى نقطة البعد الاقرب من الفلك الخارج وكذلك ايضًا بالعكْس}}}
\At{base east}{499.33pt,-194.69pt}{\hypertarget{AB01-PDF1072-L03}{}{\fontsize{13.8}{27}\ArBody\RLJ{420.10pt}{لهذه النِسْبة بالمقدار الذي به تكون الثمانية الاجزاء خمسةَ اجزاء ورُبْعًا فبه تكون الستّين جزءًا تسعةً}}}
\At{base east}{499.33pt,-221.39pt}{\hypertarget{AB01-PDF1072-L04}{}{\fontsize{13.8}{27}\ArBody\RLJ{420.10pt}{وثلثين جزءًا واثنتين وعشرين فإِذ قد وضَح مقدار خطّ \Ov{ه ف} بأنّه \Ov{لط} جزءًا و\Ov{كب} دقيقة فاذا أُضِيفَ}}}
\At{base east}{499.33pt,-247.99pt}{\hypertarget{AB01-PDF1072-L05}{}{\fontsize{13.8}{27}\ArBody\RLJ{420.10pt}{الى خطّ \Ov{ه ا} الذي هو ستّون بلَغ \Ov{صط} \Ov{كب} وهو قُطْر الفلك الخارج المركز كلّه فاذا أُخِذَ نصفه كان}}}
\At{base east}{71.13pt,-246.59pt}{{\fontsize{8}{9}\selectfont 5}}
\At{base east}{499.33pt,-275.39pt}{\hypertarget{AB01-PDF1072-L06}{}{\fontsize{13.8}{27}\ArBody\RLJ{420.10pt}{نصف قطر الفلك الخارج \Ov{مط ما}\Note{2} واذ قد عُرِفَ نصف قطر فلك التدوير بحَسَب انحرافه عن الشمس}}}
\At{base east}{499.33pt,-302.39pt}{\hypertarget{AB01-PDF1072-L07}{}{\fontsize{13.8}{27}\ArBody\RLJ{420.10pt}{وعُرِفَ ما بين المركزين ونصف قطر الفلك الخارج فانّ الذي يبقى من تمام ما في ذلك من المعرفة}}}
\At{base east}{498.93pt,-328.99pt}{\hypertarget{AB01-PDF1072-L08}{}{\fontsize{13.8}{27}\ArBody\RLJ{419.70pt}{ايضاح حِساب قوس \Ov{ط ح} المرسومة في الجَدْوَل الثالث وايضاح ما يتركَّب من التعديل المُفْرَد مع\Note{3} الثاني}}}
\At{base east}{499.33pt,-355.59pt}{\hypertarget{AB01-PDF1072-L09}{}{\fontsize{13.8}{27}\ArBody\RLJ{420.10pt}{فيما بين البُعْدَيْن على الجهة التي رَسَمْتُ في الجداول وأثْبَتُّ من ذلك في الجَدْوَل الرابع والخامس امّا}}}
\At{base east}{498.23pt,-382.29pt}{\hypertarget{AB01-PDF1072-L10}{}{\fontsize{13.8}{27}\ArBody\RLJ{419.00pt}{الذي في الرابع فإِنّه اذا كانت هذه الدرجتَيْن والاربعين الدقيقة ستّين وهي المُثْبَتة في الخامس كم يكون}}}
\At{base east}{71.43pt,-381.19pt}{{\fontsize{8}{9}\selectfont 10}}
\At{base east}{499.33pt,-409.29pt}{\hypertarget{AB01-PDF1072-L11}{}{\fontsize{13.8}{27}\ArBody\RLJ{420.10pt}{مقدار ما يتركَّب منها وما نِسْبَتُهُ الى الستّين ومعرفة ذلك كما أَصِف نُخْرِج خطّ \Ov{م ه} الى علامة \Ov{ل}}}}
\At{base east}{499.33pt,-435.89pt}{\hypertarget{AB01-PDF1072-L12}{}{\fontsize{13.8}{27}\ArBody\RLJ{420.10pt}{ونَصِل نقطة \Ov{ل} بنقطة \Ov{ز} فيصير مثلَّث \Ov{م ل ز} متناسب الأَضْلاع معلوم الزوايا ونفرِض قوس \Ov{ا م} كما}}}
\At{base east}{499.33pt,-462.89pt}{\hypertarget{AB01-PDF1072-L13}{}{\fontsize{13.8}{27}\ArBody\RLJ{420.10pt}{فرَضها بطلميوس في هذا الباب \Ov{قك} جزءًا وهي بُعْد القمر المُضْعَف عن الشمس. ولأَنَّا جعَلْنا نِسْبة الاوتار}}}
\At{base east}{498.63pt,-489.89pt}{\hypertarget{AB01-PDF1072-L14}{}{\fontsize{13.8}{27}\ArBody\RLJ{419.40pt}{المنصَّفة الى نصف القطر تكون زاوية \Ov{ل ه ز}\Note{4} ثلثين جزءًا وزاوية \Ov{ز ه ل} لتمام الرُّبْع ستّين جزءًا بالمقدار}}}
\At{base east}{499.33pt,-516.59pt}{\hypertarget{AB01-PDF1072-L15}{}{\fontsize{13.8}{27}\ArBody\RLJ{420.10pt}{الذي به تكون الدائرة التي تستدير على مثلَّث \Ov{ز ه ل} \Ov{شس} جزءًا وايضًا فإِنّ وَتَر زاوية \Ov{ل ه ز}\Note{5} *تكون}}}
\At{base west}{508.13pt,-516.59pt}{{\fontsize{8}{9}\selectfont f. 53,v.}}
\At{base east}{71.83pt,-516.19pt}{{\fontsize{8}{9}\selectfont 15}}
\At{base east}{498.63pt,-543.59pt}{\hypertarget{AB01-PDF1072-L16}{}{\fontsize{13.8}{27}\ArBody\RLJ{419.40pt}{ايضًا ثلثين جزءًا ووتر زاوية \Ov{ز ه ل} تكون قريبًا من \Ov{نا نح}\Note{6} بالمقدار الذي به يكون\Note{7} خطّ \Ov{ه ز} ستّين وامّا}}}
\At{base east}{499.33pt,-569.09pt}{\hypertarget{AB01-PDF1072-L17}{}{\fontsize{13.8}{27}\ArBody\RLJ{420.10pt}{بالمقدار الذي به يكون خطّ \Ov{ه ز} \Ov{ي} \Ov{يط} فبه يكون خطّ \Ov{ه ل} قريبًا من \Ov{ه} \Ov{ي} وخطّ \Ov{ز ل} يكون \Ov{ط} \Ov{يو}}}}
\At{base east}{499.33pt,-596.79pt}{\hypertarget{AB01-PDF1072-L18}{}{\fontsize{13.8}{27}\ArBody\RLJ{420.10pt}{واذا كان ايضًا في الشَّكْل خطّ \Ov{ه ك ن} مُمَاسًّا لفلك التدوير وموضع القمر من فلك التدوير نقطة \Ov{ك}}}}
\At{base east}{499.33pt,-623.49pt}{\hypertarget{AB01-PDF1072-L19}{}{\fontsize{13.8}{27}\ArBody\RLJ{420.10pt}{يكون اكثر ما يجتمع من الاختلاف الاوّل مع ما يتركَّب معه من الاختلاف الثاني واذ خطّ \Ov{م ك} نصف}}}
\At{base east}{499.33pt,-650.79pt}{\hypertarget{AB01-PDF1072-L20}{}{\fontsize{13.8}{27}\ArBody\RLJ{420.10pt}{قُطْر فلك التدوير وخطّ \Ov{ز م} نصف قطر الفلك الخارج وبه علِمْتَ هذه الأقدار فمن نسبة \Ov{ز م} و\Ov{ز ل}}}}
\At{base east}{71.83pt,-650.79pt}{{\fontsize{8}{9}\selectfont 20}}
\At{base east}{499.33pt,-677.09pt}{\hypertarget{AB01-PDF1072-L21}{}{\fontsize{13.8}{27}\ArBody\RLJ{420.10pt}{تكون نسبة \Ov{ل م} معلومة ولذلك يكون \Ov{م ل} كلّه \Ov{مح نج} فإِذا أُلْقِيَ منه خطّ \Ov{ه ل} الذي قد ظهَر انّه \Ov{ه} \Ov{ي}}}}
\At{west}{78.33pt,-706.49pt}{\rule{100.8pt}{0.5pt}}
\At{base west}{107.53pt,-730.38pt}{\hypertarget{AB01-PDF1072-N01}{}{\fontsize{9.2}{11}\selectfont \makebox[393.80pt][s]{1) Cod. sine articulo. — 2) Ex Platonis versione, excidisse patet \ArRun{فيكون ما بين المركزين ي يط بالمقدار}}}}
\At{base west}{81.23pt,-745.49pt}{\hypertarget{AB01-PDF1072-N02}{}{\fontsize{9.2}{11}\selectfont \makebox[420.10pt][s]{\ArRun{الذي به يكون خطّ اه ستّين جزءًا} — 3) Cod. \ArRun{من} — 4) Cod., ut Plato, \ArRun{ا ه م} — 5) Cod. \ArRun{ا ه م} — 6) Cod.}}}
\At{base west}{81.23pt,-760.59pt}{\hypertarget{AB01-PDF1072-N03}{}{\fontsize{9.2}{11}\selectfont \ArRun{يا يج} — 7) Deest in codice.}}
\DSource{1071}{PAG. 81}
\At{base east}{547.53pt,-140.15pt}{\hypertarget{AB01-PDF1071-L01}{}{\fontsize{13.8}{27}\ArBody\RLJ{420.80pt}{بقي خطّ \Ov{ه م} الذي من المركز \Ov{مج} \Ov{مج} وخطّ \Ov{م ك} الذي هو نصف قطر فلك التدوير فقد كان ظهَر}}}
\At{base east}{548.63pt,-166.75pt}{\hypertarget{AB01-PDF1071-L02}{}{\fontsize{13.8}{27}\ArBody\RLJ{421.90pt}{انّه \Ov{ه يه} ولكن بالمقدار الذي به يكون خطّ \Ov{ه م} الذي من المركز ستّين فبه يكون خطّ \Ov{ك م} وهو نصف}}}
\At{base east}{549.33pt,-193.35pt}{\hypertarget{AB01-PDF1071-L03}{}{\fontsize{13.8}{27}\ArBody\RLJ{422.60pt}{قطر فلك التدوير المنحرِف سبعة اجزاء واثنتيْ عشرة دقيقة بالتقريب والقوس التي عليه قريبة من}}}
\At{base east}{549.33pt,-219.35pt}{\hypertarget{AB01-PDF1071-L04}{}{\fontsize{13.8}{27}\ArBody\RLJ{422.60pt}{ستّة\Note{1} اجزاء واربع وخمسين دقيقة وذلك هو مقدار قوس \Ov{م ك} فإِذا طُرِحَ من ذلك الخمسة الاجزاء}}}
\At{base east}{549.73pt,-246.35pt}{\hypertarget{AB01-PDF1071-L05}{}{\fontsize{13.8}{27}\ArBody\RLJ{423.00pt}{والدقيقة التي هي قدر الاختلاف المفرد كلّه بقي الذي يتركَّب معه من الاختلاف الثاني جزءًا وثلثًا}}}
\At{base west}{556.03pt,-247.05pt}{{\fontsize{8}{9}\selectfont 5}}
\At{base east}{550.13pt,-272.95pt}{\hypertarget{AB01-PDF1071-L06}{}{\fontsize{13.8}{27}\ArBody\RLJ{423.40pt}{وخمسين دقيقة واذا كانت الدرجتانِ والثُّلْثا درجة بالتقريب ستّين كانت هذه الدرجة والثلث والخمسون\Note{2}}}}
\At{base east}{550.43pt,-299.55pt}{\hypertarget{AB01-PDF1071-L07}{}{\fontsize{13.8}{27}\ArBody\RLJ{423.70pt}{دقيقة من الستّين \Ov{مب لح} وهي المرسومة تحت \Ov{قك} في الجدول الرابع وقد حصَلت بنسبة الدقائق الى}}}
\At{base east}{550.13pt,-326.25pt}{\hypertarget{AB01-PDF1071-L08}{}{\fontsize{13.8}{27}\ArBody\RLJ{423.40pt}{الجزء الواحد وهي نسبة \Ov{\AbjadZero{} مب} \Ov{لح} الى ستّين واذا مُدَّت هذه الاثنان والاربعون والثماني والثلثون فجُعِلَتْ}}}
\At{base east}{550.43pt,-352.85pt}{\hypertarget{AB01-PDF1071-L09}{}{\fontsize{13.8}{27}\ArBody\RLJ{423.70pt}{ستّين كان الجزء والثلثة والخمسون\Note{2} دقيقة حينئذ *جزءَيْن وتسعًا وثلثين دقيقة وهي المرسومة في الخامس}}}
\At{base east}{118.93pt,-352.85pt}{{\fontsize{8}{9}\selectfont f. 54,r.}}
\At{base east}{550.43pt,-379.55pt}{\hypertarget{AB01-PDF1071-L10}{}{\fontsize{13.8}{27}\ArBody\RLJ{423.70pt}{تحت \Ov{قك}. وايضًا نعلَم ما بين البعد الابعد الحقّيّ والبعد الاوسط الذي هو قوس \Ov{ط ح} على الرسم}}}
\At{base west}{558.13pt,-379.85pt}{{\fontsize{8}{9}\selectfont 10}}
\At{base east}{550.43pt,-406.15pt}{\hypertarget{AB01-PDF1071-L11}{}{\fontsize{13.8}{27}\ArBody\RLJ{423.70pt}{الموضوع بأن نفرِض بُعْد القمر عن الشمس بمسيرها الاوسط المُضْعَف \Ov{ص ل} كما جعَله بطلميوس في الشكل}}}
\At{base east}{550.43pt,-432.75pt}{\hypertarget{AB01-PDF1071-L12}{}{\fontsize{13.8}{27}\ArBody\RLJ{423.70pt}{الذي انتهت اليه دَلالته على ما ذكرنا لتكون حركة القمر في فلك تدويره من نقطة \Ov{ط} \Ov{شلج} \Ov{يب}}}}
\At{base east}{550.43pt,-459.45pt}{\hypertarget{AB01-PDF1071-L13}{}{\fontsize{13.8}{27}\ArBody\RLN{136.60pt}{ونرسُم هذه الدائرة مِثالًا لذلك.}}}
\begin{tikzpicture}[remember picture,overlay]\draw[line width=1.47pt] ([shift={(223.94pt,-576.00pt)}]current page.north west) circle[radius=71.70pt];\draw[line width=1.37pt] ([shift={(157.77pt,-604.30pt)}]current page.north west) circle[radius=23.81pt];\draw[line width=1.32pt] ([shift={(223.86pt,-504.30pt)}]current page.north west) -- ([shift={(224.18pt,-647.70pt)}]current page.north west);\draw[line width=1.43pt] ([shift={(224.02pt,-576.31pt)}]current page.north west) -- ([shift={(158.72pt,-604.59pt)}]current page.north west);\draw[line width=1.34pt] ([shift={(231.50pt,-592.40pt)}]current page.north west) -- ([shift={(134.33pt,-608.48pt)}]current page.north west);\draw[line width=1.33pt] ([shift={(134.05pt,-602.26pt)}]current page.north west) -- ([shift={(224.10pt,-610.73pt)}]current page.north west);\draw[line width=1.30pt] ([shift={(224.02pt,-577.47pt)}]current page.north west) -- ([shift={(231.50pt,-592.40pt)}]current page.north west);\draw[line width=1.12pt] ([shift={(217.95pt,-594.64pt)}]current page.north west) -- ([shift={(224.10pt,-609.76pt)}]current page.north west);\draw[line width=1.34pt] ([shift={(158.34pt,-604.56pt)}]current page.north west) -- ([shift={(142.92pt,-622.90pt)}]current page.north west);\end{tikzpicture}
\At{center}{223.47pt,-488.45pt}{\hypertarget{AB01-PDF1071-F01-T01}{}{\fontsize{15.5}{18.6}\ArBody\RL{ا}}}
\At{center}{229.29pt,-573.76pt}{\hypertarget{AB01-PDF1071-F01-T02}{}{\fontsize{15.5}{18.6}\ArBody\RL{د}}}
\At{center}{237.39pt,-587.81pt}{\hypertarget{AB01-PDF1071-F01-T03}{}{\fontsize{15.5}{18.6}\ArBody\RL{ك}}}
\At{center}{216.69pt,-590.42pt}{\hypertarget{AB01-PDF1071-F01-T04}{}{\fontsize{15.5}{18.6}\ArBody\RL{س}}}
\At{center}{227.97pt,-596.57pt}{\hypertarget{AB01-PDF1071-F01-T05}{}{\fontsize{15.5}{18.6}\ArBody\RL{ه}}}
\At{center}{218.31pt,-617.45pt}{\hypertarget{AB01-PDF1071-F01-T06}{}{\fontsize{15.5}{18.6}\ArBody\RL{ز}}}
\At{center}{157.11pt,-612.17pt}{\hypertarget{AB01-PDF1071-F01-T07}{}{\fontsize{15.5}{18.6}\ArBody\RL{ب}}}
\At{center}{128.85pt,-594.83pt}{\hypertarget{AB01-PDF1071-F01-T08}{}{\fontsize{15.5}{18.6}\ArBody\RL{ط}}}
\At{center}{129.93pt,-614.81pt}{\hypertarget{AB01-PDF1071-F01-T09}{}{\fontsize{15.5}{18.6}\ArBody\RL{ح}}}
\At{center}{141.21pt,-631.31pt}{\hypertarget{AB01-PDF1071-F01-T10}{}{\fontsize{15.5}{18.6}\ArBody\RL{م}}}
\At{center}{225.51pt,-667.25pt}{\hypertarget{AB01-PDF1071-F01-T11}{}{\fontsize{15.5}{18.6}\ArBody\RL{ج}}}
\At{base east}{526.73pt,-486.05pt}{\hypertarget{AB01-PDF1071-L14}{}{\fontsize{13.8}{27}\ArBody\RLN{215.30pt}{قال هذه دائرة الفلك الخارج عليها \Ov{ا ب ج} على مركز}}}
\At{base east}{550.43pt,-512.35pt}{\hypertarget{AB01-PDF1071-L15}{}{\fontsize{13.8}{27}\ArBody\RLN{237.20pt}{\Ov{د}\Note{3} وقُطْر \Ov{ا ج} وعليه يقَع مركز فلك البروج في موضع علامة}}}
\At{base west}{559.23pt,-513.45pt}{{\fontsize{8}{9}\selectfont 15}}
\At{base east}{550.43pt,-539.35pt}{\hypertarget{AB01-PDF1071-L16}{}{\fontsize{13.8}{27}\ArBody\RLN{239.40pt}{\Ov{ه}\Note{4} والدائرة على نقطة \Ov{ب} هي المسمَّاة بفلك التدوير}}}
\At{base east}{550.43pt,-564.55pt}{\hypertarget{AB01-PDF1071-L17}{}{\fontsize{13.8}{27}\ArBody\RLN{237.60pt}{عليها \Ov{م ح ط} ونُخْرِج خطّ \Ov{ب م} وخطّ \Ov{ه ب ح} ونَصِل خطّ}}}
\At{base east}{550.43pt,-592.25pt}{\hypertarget{AB01-PDF1071-L18}{}{\fontsize{13.8}{27}\ArBody\RLN{237.20pt}{\Ov{ب ه} بنقطة \Ov{ك} ونَصِل نقطة \Ov{ك} بنقطة \Ov{د} فتكون زاوية}}}
\At{base east}{550.43pt,-619.25pt}{\hypertarget{AB01-PDF1071-L19}{}{\fontsize{13.8}{27}\ArBody\RLN{236.80pt}{\Ov{ك د ه} من المثلَّث النصف الجزء الزائد على \Ov{ص} وقوس}}}
\At{base east}{550.43pt,-645.95pt}{\hypertarget{AB01-PDF1071-L20}{}{\fontsize{13.8}{27}\ArBody\RLN{236.80pt}{\Ov{ه ك} نصف جزء بالمقدار الذي به تكون الدائرة المستديرة}}}
\At{base west}{559.63pt,-646.25pt}{{\fontsize{8}{9}\selectfont 20}}
\At{base east}{550.43pt,-671.85pt}{\hypertarget{AB01-PDF1071-L21}{}{\fontsize{13.8}{27}\ArBody\RLN{236.80pt}{على مثلَّث \Ov{د ك ه} \Ov{شس} جزءًا ووَتَرها المنصَّف \Ov{\AbjadZero{} لا كه} بالمقدار}}}
\At{base east}{550.13pt,-699.14pt}{\hypertarget{AB01-PDF1071-L22}{}{\fontsize{13.8}{27}\ArBody\RLJ{423.40pt}{الذي به يكون خطّ \Ov{ه د} الذي هو نصف القطر ستّين وتبقى زاوية \Ov{ك ه د} *\Ov{فط ل} وقوس \Ov{ك د} لذلك}}}
\At{base east}{119.73pt,-699.14pt}{{\fontsize{8}{9}\selectfont f. 54,v.}}
\At{base west}{513.13pt,-775.14pt}{{\fontsize{9}{10}\selectfont\bfseries 11}}
\At{west}{126.93pt,-726.89pt}{\rule{101.1pt}{0.5pt}}
\At{base west}{155.33pt,-754.95pt}{\hypertarget{AB01-PDF1071-N01}{}{\fontsize{9.2}{11}\selectfont 1) Cod. \ArRun{خمسة} — 2) Cod. sine articulo. — 3) Cod. \ArRun{ه} — 4) Cod. \ArRun{د}}}
\end{document}
